\subsection{A prescribed partial model can force descent to form degree zero}

The next two statements concern the particular partial model and its
particular maps.  They give an integral admissible target and a choice
of the actual construction in SA08, C21, for which the truncation
agreements required in C34 cannot have depths tending to infinity.
The target will also have an explicit section and retraction
$G\to A\to G$.  These maps specify the scope of the obstruction.

\begin{proposition}[AT11: Proved; arbitrarily long coefficient-syzygy descent]
\label{AT11}
Let $p$ be a prime, let $R$ be either
$\mathbb Z[1/(p-1)!]$ or $\mathbb Z_{(p)}$, and put
$S_R=\Gamma_R(T)$ and $J=(p,I_1)\subset S_R$.
There are locally finite partial graded free resolutions of arbitrary
specified finite length
\[
 0\longrightarrow K_L\longrightarrow F_{L-1}
 \longrightarrow\cdots\longrightarrow F_0
 \xrightarrow{\pi_J}J\longrightarrow0,\qquad L\geq1,
 \tag{AT.30}
\]
whose free generators have nonnegative weights, with
\[
 F_0=S_Ru\oplus\bigoplus_{a\geq1}S_Re_a,
 \quad \operatorname{wt}(u)=0,\quad\operatorname{wt}(e_a)=a,
 \quad \pi_J(u)=p,\quad\pi_J(e_a)=T^{[a]}.
 \tag{AT.31}
\]
Every weight group of every kernel and displayed module is
finite free over $R$.
Here ``locally finite'' means that finitely many free generators
occur in every bounded interval of weights.

Put $D=S_R/pS_R=\Gamma_{\mathbb F_p}(T)$ and use a bar for
reduction modulo $p$.  For every $M\geq1$, pulling back
\textup{(AT.30)} modulo $p$ along the inclusion
$I_M(D)\to\overline J$ gives a nonzero class
\[
 \chi_{M,L}\in
 \operatorname{Ext}^L_{\mathrm{gr}D}(I_M(D),\overline K_L).
 \tag{AT.32}
\]
In particular these classes persist under every further coefficient
tail restriction, for every chosen length $L$.
\end{proposition}

\begin{proof}
Every weight of $F_0$ is finite free over the PID $R$, as is
every weight of $J$.  Let $K_1=\ker\pi_J$.
Inductively choose an $R$-basis in every weight of $K_i$ and
let $F_i$ be free over $S_R$ on the union of these bases;
map each generator to the basis element it names and put
$K_{i+1}=\ker(F_i\to K_i)$.  At every finite induction step
the modules have nonnegative weights and finite free weight
groups: in weight $w$ only the finitely many chosen generators
of weights at most $w$ occur.  The kernel is a submodule of
a finite free module over a PID and is therefore finite free.
This proves every assertion about \textup{(AT.30)}.  Each short
exact sequence has $R$-free weightwise cokernel, so reduction
modulo $p$ preserves its exactness.

There is an explicit decomposition of graded $D$-modules
\[
 \overline J=\mathbb F_p\,\overline p\ \oplus\ I_1(D).
 \tag{AT.33}
\]
Its weight-zero generator is the class of $p$ in $pR/p^2R$.
Every positive-weight coefficient kills this generator since
$pT^{[a]}\in pJ$ for $a>0$.  In positive weight, $J^a=R T^{[a]}$
and $(pJ)^a=pR T^{[a]}$, giving exactly the second summand.
The map $I_M(D)\to\overline J$ in the assertion lands in
that second summand.

This map has no graded $D$-linear lift through $\overline F_0$.
Indeed choose $m=p^r\geq M$.  A homogeneous putative lift has
the value
\[
 T^{[m]}\longmapsto
 c_uT^{[m]}u+\sum_{a=1}^m c_aT^{[m-a]}e_a.
\]
The characteristic-$p$ identity $(1+X)^{p^r}=1+X^{p^r}$
gives $\binom ma=0$ for $0<a<m$.  The $u$ term maps to zero
in the positive-weight part of \textup{(AT.33)}.
Consequently the lifting equation forces $c_m=1$.
Multiplication by $T^{[(p-1)m]}$ kills $T^{[m]}$ in $I_M(D)$,
since $\binom{pm}{m}=0$ in $\mathbb F_p$.
It does not kill the $e_m$ component of its asserted lift:
that component becomes $T^{[(p-1)m]}e_m\ne0$ in a free
summand of $\overline F_0$.  The other free summands cannot
cancel it.  This proves the nonlifting assertion for every $M$.

We explain precisely why it detects every length, including
the infinite-rank issue.  The graded module $D$ is injective
over itself by Nagpal--Snowden, Theorem~7.1(a), printed
p.~39, with proof on p.~41 of arXiv:1606.03431v2
\cite{IC-nagpal-snowden}.  We use the full integer-graded
module category and weight-preserving maps.  Each shift of
$D$ is injective in that category.  For each $\overline F_i$
the direct sum of its free shifted summands is the graded
product of those same summands: in any weight, only finitely
many summands are nonzero.  Products of injectives are
injective, because Hom into the product is the product of
the exact Hom functors.  Thus every $\overline F_i$ is
injective in this specified category.

The first connecting map sends the inclusion in
\textup{(AT.33)} to a nonzero class in
$\operatorname{Ext}^1(I_M(D),\overline K_1)$, exactly by
the proved failure to lift.  For $i\geq2$, the connecting
map associated to
$0\to\overline K_i\to\overline F_{i-1}
\to\overline K_{i-1}\to0$ is an isomorphism
\[
 \operatorname{Ext}^{i-1}(I_M(D),\overline K_{i-1})
 \xrightarrow{\ \cong\ }
 \operatorname{Ext}^{i}(I_M(D),\overline K_i),
\]
since the two intervening positive Ext groups into
$\overline F_{i-1}$ vanish.  Their composite is the
Yoneda class of the pullback of \textup{(AT.30)}.
It is nonzero for every $M$, proving \textup{(AT.32)}
and persistence.
\end{proof}

Here is the explicit cochain consequence of AT11 that we
will use.  Choose integers $q>L\geq1$ and a weight shift $W$.
Form a cochain complex $C$ by putting $K_L\langle W\rangle$
in form degree $q-L-1$, $F_i\langle W\rangle$ in degree
$q-1-i$, and $S_R\langle W\rangle\vartheta$ in degree $q$.
The last differential is
$F_0\langle W\rangle\xrightarrow{\pi_J}
J\langle W\rangle\hookrightarrow S_R\langle W\rangle$;
the other differentials are those of \textup{(AT.30)}.
Let $Q_\bullet\to I_M\langle W\rangle$ be a graded free
resolution whose reduction is still a resolution; the
weightwise construction in the preceding proof supplies one.
The closed top map
$I_M\langle W\rangle\to C^q$, $a\mapsto a\vartheta$,
has nullhomotopies on this resolution, but every such
nullhomotopy has a nonzero component
$Q_L\to K_L\langle W\rangle$.

To prove the last sentence, lift the inclusion $I_M\to J$
successively along the free resolution using projectivity
of the $Q_i$.  At the last step factor the composite into
the kernel $K_L$.  This constructs a nullhomotopy, with
signs adjusted successively by $(-1)^i$ if required by the
Hom differential convention.  Its terminal equation
$h_Ld_Q=0$ follows by applying the injective map
$K_L\hookrightarrow F_{L-1}$ and using the preceding
lifting equation together with $d_Q^2=0$.
Conversely, write any such
nullhomotopy as maps
$h_i:Q_i\to C^{q-1-i}$.  Its degree-zero equation says
that the differential of $h_0$ is the given inclusion
$I_M\langle W\rangle\to S_R\langle W\rangle$.
Before reducing modulo $p$, factor this equation through
the monomorphism $J\hookrightarrow S_R$.  It says exactly
that $\pi_J h_0$ is the prescribed inclusion into $J$.
The remaining equations are successive lifts of the
syzygies.  After reduction modulo $p$, the cocycle
$Q_L/pQ_L\to\overline K_L\langle W\rangle$ so obtained
represents \textup{(AT.32)}.  If $h_L$ were zero this
class would be zero, a contradiction.  The factorization
before reduction is essential: the map
$\overline J\to D$ kills the first summand of
\textup{(AT.33)}, whereas $J\to S_R$ is injective.

\begin{theorem}[AT12: Proved; one prescribed C21 model has no growing local completion depth]
\label{AT12}
There is an integral admissible strict PD theory $A$, with
flat terms and support $q\leq w$, and there is a choice of
the simultaneous C21 construction giving actual maps
\[
                 A\longleftarrow P\longrightarrow G,
 \tag{AT.34}
\]
with all of C21's stated properties and $P_0=S$.
For every odd prime $p$, there is no admissible theory $E$
over $\mathbb Z_{(p)}$, map $g:E\to A_{(p)}$, and
isomorphism of truncated simplicial PD-DGAs
\[
 \alpha:(P_{(p)})_{<2}\xrightarrow{\ \cong\ }E_{<2}
 \quad\hbox{with}\quad
 g_{<2}\alpha=(P\to A)_{(p),<2}.
 \tag{AT.35}
\]
Consequently the specified map
$P\to A\times_{\mathbb Z}G$ has no family of the local
truncation agreements of C34 with $d_p\to\infty$.
The fiber product here uses the augmentations to the
weight-zero coefficient ring.  There are also actual
strict PD maps $G\to A\to G$ with composite the identity.
\end{theorem}

\begin{proof}
\emph{The primewise algebra.}
For each odd prime $p$, put
\[
 R_p=\mathbb Z[1/(p-1)!],\qquad b_p=p+1,
 \qquad q_p=2p,\qquad W_p=p b_p,\qquad L_p=2p-1.
\]
Initially work over $S_{R_p}$.  Start with
$\Lambda(x)\otimes\Gamma(y)$, $|x|=1$, $|y|=2$,
$dx=y$, giving both generators weight $b_p$.
Quotient by all monomials of total arity greater than $p$.
This is a differential PD ideal: differentiation preserves
arity, multiplication adds arities, and an $n$th divided
power multiplies positive arity by $n$.
For $1\leq m<p$, its arity-$m$ pair is
\[
 x\gamma_{m-1}(y)\xrightarrow{\ m\ }\gamma_m(y),
 \tag{AT.36}
\]
and contracts over $R_p$.  Write
$u=x\gamma_{p-1}(y)$ and $\vartheta=\gamma_p(y)$.
They have weight $W_p$, form degrees $q_p-1,q_p$,
and $du=p\vartheta$.  Both annihilate the positive
arity ideal.

Replace this last pair by the complex $C$ following
AT11, with $q=q_p$, $L=L_p$, and shift $W=W_p$.
Identify its named generator $u\in F_0\langle W_p\rangle$
with the old product $x\gamma_{p-1}(y)$.
Every new summand annihilates every positive arity old
element and every new summand.  Coefficients act by their
given $S_{R_p}$-module action.  Retain all old products,
and put all higher divided powers of new even elements
equal to zero.  The resulting augmented differential
graded algebra is denoted $B_p$, and its augmentation
ideal over $S_{R_p}$ is $B_{p,+}$.

Here are full checks of the additional structure.  The
old top elements already annihilate positive arity, so
identifying $u$ introduces no inconsistent product.
The image of the differential of a new element stays
in the new modules or in the span of $u,\vartheta$;
all these elements annihilate positive arity.  The
Leibniz identity on a product involving a new element
therefore has both sides zero unless the other factor
is a coefficient, where it is $S_{R_p}$-linearity.
On old products it is the identity in the truncated
free disk, including $du=p\vartheta$.
Every weight group is torsionfree over $R_p$.
Embed the constructed algebra into its rationalization
and use $\gamma_n(a)=a^n/n!$ on even positive-weight
elements.  These operations preserve the integral
lattice: on the old disk they are its original divided
powers; on the square-zero new summands they vanish for
$n>1$; and on a sum the formula is
$\gamma_n(a+v)=\gamma_n(a)+\gamma_{n-1}(a)v$.
If $a$ contains an old positive arity component its
product with $v$ is zero, leaving only the allowed
coefficient action.  The same argument handles finite
sums of new elements.  Thus all divided-power identities
and the differential identity hold after the injective
rationalization and hence over $R_p$ itself.
This proves the claimed strict PD structure.

The arity decomposition of its underlying cochain
module is its coefficient copy, the pairs
\textup{(AT.36)}, and the specified complex $C$.
The pairs contract, and \textup{(AT.30)} gives
\[
 H^0(B_p)=S_{R_p},\qquad
 H^{q_p}(B_{p,+})=(S_{R_p}/J)\langle W_p\rangle
                   =\mathbb F_p\langle W_p\rangle,
 \tag{AT.37}
\]
with all other cohomology of $B_{p,+}$ zero.
Its bottom new module $K_{L_p}$ is in form degree zero
and injects into the next term; it adds no $H^0$.
Every homogeneous term is finite free over $R_p$.
All old arity-$m$ terms satisfy $q\leq2m\leq m b_p$;
all new terms have weight at least $W_p\geq q_p$.
This verifies support $q\leq w$.

\emph{One integral algebra.}
Put
\[
 B=S\oplus\bigoplus_{p\text{ odd}}B_{p,+}.
 \tag{AT.38}
\]
The coefficient action on the $p$th ideal is induced
by $S\to S_{R_p}$, and products between two different
positive ideals are zero.  More explicitly, this
algebra embeds into $\prod_{p\text{ odd}}B_p$ by
$(s,(v_p))\mapsto(s+v_p)_p$.
It is closed under every operation: applying an
operation to finitely many elements changes the common
coefficient part in the same way at every coordinate
and changes only their finitely many supported ideals.
For divided powers this follows from the PD sum formula.
Thus \textup{(AT.38)} is an integral strict PD differential
algebra, with coefficient section and augmentation.
Its terms are flat over $\mathbb Z$, being sums of
free $R_p$-modules.  In each fixed weight only finitely
many primes occur, since $b_p\to\infty$.
Its extra cohomology is exactly that in
\textup{(AT.37)}, one class at each indicated
form degree and weight.

\emph{The simplicial parameter and admissible theory.}
Let
\[
 C_n=\mathbb Z^{\{(i,j,k):0\leq i\leq j\leq k\leq n\}},
 \tag{AT.39}
\]
with coordinatewise multiplication.  For a monotone
map $\eta:[m]\to[n]$ the simplicial structure map is
$f\mapsto((i,j,k)\mapsto f(\eta(i),\eta(j),\eta(k)))$.
Use normalization $N_n=\bigcap_{i=1}^n\ker d_i$ and
differential $d_0$.  Its basis consists precisely of
indicators of triples containing every vertex
$1,\ldots,n$.  The full normalized complex is
\[
\begin{array}{c|c}
 n&\text{basis of }N_nC\\\hline
0&\epsilon_{000}\\
1&\epsilon_{001},\epsilon_{011},\epsilon_{111}\\
2&\epsilon_{012},\epsilon_{112},\epsilon_{122}\\
3&\epsilon_{123}
\end{array}
\]
and is zero in higher degrees.  The only nonzero
differentials on this basis are
$\epsilon_{111}\mapsto\epsilon_{000}$,
$\epsilon_{112}\mapsto\epsilon_{001}$,
$\epsilon_{122}\mapsto\epsilon_{011}$,
$\epsilon_{123}\mapsto\epsilon_{012}$.
Sending each listed image to its unique listed
preimage and the remaining basis elements to zero
is a contraction.  In particular $\epsilon_{012}$
is a closed normalized degree-two element.

Take the standard theory $G$ defined in
\textup{(AL.4)}.  Its terms are free over $S$, its
simplexwise coefficient map has a weight-preserving
$S$-linear vertical contraction, and its rows are
contractible for $w>q$, as proved there.
Let $R_n'=G_n\otimes_S B$ and define
\[
 A_n=\mathop{\prod\nolimits_{G_n}}_{0\leq i\leq j\leq k\leq n}
       R_n'
   =G_n\oplus(G_n\otimes_S B_+)\otimes_{\mathbb Z}C_n.
 \tag{AT.40}
\]
The first expression is a fiber product of the
augmentations $R_n'\to G_n$.
It defines every algebra operation coordinatewise.
Faces and degeneracies apply the given map of $G$
and the coordinate restriction in \textup{(AT.39)}.
They are strict PD maps and satisfy the simplicial
identities by composition of the monotone maps.
The ordinary tensor product here has its strict PD
structure: in simplex $n$, a choice of a complement
to the marked coefficient in the defining module of
$G_n$ writes it as $S$ tensored with free even-to-odd
PD disks.  Tensoring these disks with $B$ defines
the structure, and the maps induced by $G$ preserve
it by their original PD identities.  This also
proves flatness of all terms of \textup{(AT.40)}.
The coefficient section and augmentation give the
asserted maps $G\to A\to G$.

Tensor the $S$-linear vertical contraction of $G_n$
with $B_+$ and the finite free group $C_n$.
It computes all additional simplexwise cohomology
as $C_n\otimes\mathbb F_p$ in form degree $2p$
and weight $W_p$, and none elsewhere.
On the other hand each added horizontal form-weight
row is contractible: tensor the additive simplicial
contraction of $C$ with the other simplicial factor,
using the corresponding degeneracy on that factor
in each homotopy operator.  This proves, uniformly
in the simplex, the admissibility bounds
\[
 \mu_A(d)=d+1,\qquad
 \lambda_A(d)=1+\max\bigl(\{W_p: p\text{ odd},\ 2p\leq d\}
                                      \cup\{0\}\bigr).
 \tag{AT.41}
\]
Products and tensor factors preserve $q\leq w$,
so $A$ has the required support as well.

\emph{A legitimate choice of the C21 partial model.}
Run the construction of C21 simultaneously for
$A$ and $G$, choosing $W(q)=q+1$ for its weight
thresholds.  At its form-zero stage, take all
maps to $A$ through the displayed section $G\to A$,
and choose the same primitives as for the map to
$G$.  This is possible because the required rows
of $G$ are contractible.  It ensures that the image
of every final form-zero element of $P$ lies in
$G\subset A$: all later stages add generators only
in form degrees at least one.

At the first form-one round, for each odd $p$
consider the zero normalized cycle in homological
degree one and weight $b_p$.  C21 attaches a cell
for this cycle, with parameter inclusion
$K(1)\subset D(2)$.
Send the bottom generator to zero in both targets,
the top generator to zero in $G$, and the top
generator to
\[
       x_p\epsilon_{012}\in N_2A^{b_p,1}
 \tag{AT.42}
\]
in $A$.  It is a valid parameter map because the
element in \textup{(AT.42)} is closed for $d_0$.
Dold--Kan therefore gives an actual simplicial
linear map from $D(2)$, and the free vertical-disk
universal property gives the required strict PD
map on the cell.  The weights $b_p$ are distinct,
so these choices are simultaneous choices for
different cells of the stated construction.
All remaining choices and all further rounds
can be made by the horizontal acyclicity already
proved.  Free extensions preserve each old
element, and the construction has all of C21's
properties.  It also has $P_0=S$: in simplex zero
the inclusion $K(n)\to D(n+1)$ has zero quotient
for every $n\geq0$, so no cell adds an independent
generator there.  In particular the chosen odd
generators occur in simplex two, not in simplex zero.
Let $\xi_p\in P_2^{b_p,1}$ denote the element
whose image is \textup{(AT.42)}.

\emph{The forced contradiction to a truncation agreement.}
Fix an odd $p$ and suppose \textup{(AT.35)} is given.
Work from now on over $R=\mathbb Z_{(p)}$.
The isomorphism transports $\xi_p$ to
$\xi\in E_2^{b_p,1}$, and its prescribed image in
$A_{(p)}$ is still \textup{(AT.42)}.
There is a specified weight-preserving strict PD
cochain map
\[
 \rho:(A_{(p)})_2\longrightarrow (B_p)_{(p)}:
 \quad \text{coordinate }(0,1,2),\quad
 B\to B_p,\quad G_2\to G_0=S.
 \tag{AT.43}
\]
In this formula $B\to B_p$ extends coefficients to
$R_p$, keeps its $p$th positive ideal, and kills
the other ideals; subsequent localization is at $p$.
The last arrow is any chosen simplicial vertex
map on the $G$ factor.  Each operation in
\textup{(AT.43)} is a cochain PD map.  Their
composite is used as a map of the fixed simplex
algebras; no assertion that it is a simplicial
vertex map of $A$ is needed or made.
It sends $g(\xi)$ to $x_p$.  Moreover
\[
 \rho g(E_2^{w,0})\subset S_R^w
                    \quad\text{for every }w,
 \tag{AT.44}
\]
because the form-zero truncation is surjectively
identified with that of $P$, whose map to $A$
factors through $G$ there.

Put $q=2p$, $W=W_p$, $L=q-1$, and
$\theta_E=\gamma_p(d_E\xi)\in Z^q(E_2)^{W}$.
The compatibility of $g$ with differential and
divided powers gives $\rho g(\theta_E)=\vartheta$.
This conclusion uses no agreement about the outgoing
differential in the truncation.
Choose $M\geq1$ so large that all weights at least
$M+W$ satisfy the vertical coefficient comparison
for $E$ through form degree $q$.  Choose a graded
free resolution $Q_\bullet\to I_M\langle W\rangle$
generated entirely in weights at least $M+W$,
with finite free weight groups and kernels; the
construction in AT11, applied to this module,
does so.  Its reduction modulo $p$ stays exact.

We give the exact derived calculation that is
used here.  Over $S_R$ in nonnegative weights,
the module $S_R$ has injective dimension at most
one.  In fact, putting $S_p=\Gamma_{\mathbb Z_p}(T)$
and $E_p'=S_{\mathbb Q_p}/S_p$, the following is
an explicit injective resolution:
\[
 0\longrightarrow S_R\longrightarrow
 S_{\mathbb Q}\oplus S_{\mathbb Q_p}
 \xrightarrow{(a,b)\mapsto(b-a,\overline b)}
 S_{\mathbb Q_p}\oplus E_p'\longrightarrow0.
 \tag{AT.45}
\]
The first map is diagonal.  Its kernel assertion
is $\mathbb Q\cap\mathbb Z_p=\mathbb Z_{(p)}$
weight by weight.  Its surjectivity follows from
$\mathbb Q+\mathbb Z_p=\mathbb Q_p$: subtract the
finite negative part of a $p$-adic expansion to
obtain the rational summand, and then lift an
arbitrary class in $E_p'$.  Rational injectivity
is exactly the colimit formula \textup{(AT.7)}
over $R$.  The injectivity of $E_p'$ over $S_p$
is proved in AT06, and it remains injective on
restriction to $S_R$ by adjunction along the flat
extension $S_R\to S_p$.  This proves
\textup{(AT.45)} and the asserted dimension bound.

Consider the total Hom complex
$\operatorname{Hom}_{\mathrm{gr}S_R}(Q_\bullet,E_2)$.
Its cochain degree $a$ is the finite sum of
$\operatorname{Hom}(Q_i,E_2^{a-i})$, $0\leq i\leq a$.
The first-quadrant hypercohomology spectral sequence
has terms
\[
 E_2^{i,j}=\operatorname{Ext}^i_{\mathrm{gr}S_R}
           (I_M\langle W\rangle,H^j(E_2))
 \quad\Longrightarrow\quad
 H^{i+j}\operatorname{Hom}(Q_\bullet,E_2).
 \tag{AT.46}
\]
For $1\leq j\leq q$, the Hom complex from $Q_\bullet$
to $H^j(E_2)$ is identically zero: every free
generator of $Q$ is above the chosen vertical
cutoff.  For $j=0$ that Hom complex is identified,
by the coefficient map in those weights, with
$\operatorname{Hom}(Q_\bullet,S_R)$.
Thus every $E_2$ term of total degree $q$ is zero,
the last one by $q>1$ and \textup{(AT.45)}.
Every relevant diagonal is finite, so this gives
$H^q\operatorname{Hom}(Q_\bullet,E_2)=0$ without
an unbounded convergence assertion.

The closed map
$I_M\langle W\rangle\to E_2^q$, $a\mapsto a\theta_E$,
therefore has a nullhomotopy on $Q_\bullet$.
Compose it with $\rho g$ and then with the cochain
projection of $(B_p)_{(p)}$ onto its direct summand
$C$ used in constructing the top arity.
The resulting nullhomotopy is for
$a\mapsto a\vartheta$.  Its component into $C^0$
is zero by \textup{(AT.44)}, since the projection
to $C$ kills the coefficient copy of $S_R$.
The cochain consequence of AT11 with $L=q-1$
asserts that precisely this bottom component must
be nonzero.  This contradiction proves
\textup{(AT.35)} impossible.

Finally any C34 agreement of depth $d_p\geq2$
for the specified map to $A\times_{\mathbb Z}G$
restricts to \textup{(AT.35)} after the actual
projection of that fiber product to $A$.
Hence all its odd-prime depths must be at most one,
and cannot tend to infinity.  The displayed maps
$G\to A\to G$ have already been constructed, so
this failure for the prescribed diagram does not
produce a pair without a common admissible source.
\end{proof}
